% Disjoint continuation; core, localization and BVM fragments remain frozen.
\subsubsection{The full-datum closed-base construction and collapse}

The original p.~332 theorem distinguishes a closed-category base from a
geometric base \cite[p.~332]{Lawvere1970}; the event is Nice 1970 and the
proceedings publication is 1971. The following construction retains the
given finite-limit functor \(u_*:\mathcal E\to\mathcal E_0\) and
Heyting homomorphism \(h:u_*T\to T_0\). It does not assume that
\(u_*\) itself descends through the selected-mono localization.

For the historical terminology, a closed functor has a functor, a unit
map, and a natural internal-hom comparison satisfying identity, unit and
composition compatibilities. The hom comparison need not be invertible.
Eilenberg--Kelly distinguish this from \emph{normality}, which identifies
the underlying-element functors \cite[pp.~433--434,439]{EilenbergKelly1966}.
Their relative-base paragraph imposes a stronger normal convention
\cite[p.~549]{EilenbergKelly1966}. Lawvere's six-page article does not
print such an additional element-functor requirement. Here we construct
the actual closed functor, its base-valued hom objects and comparison;
we do not silently identify the resulting enriched global underlying
category with the ordinary localization. The countercase below explains
why that stronger identification cannot follow from the full printed datum.

\begin{hbdefinition}{def:p10-lawvere-closed-base-partials}{The h-total partial-element quotient}
Let \(\partial X\subseteq\mathcal P_{\mathcal E}X\) be the object
of subsingleton predicates, with universal partial member
\(D_X\hookrightarrow\partial X\) and evaluation \(D_X\to X\).
Write \(d_X:\partial X\to T\) for its domain classifier. On
\(\partial X\times\partial X\), let
\(e_X(p,q)\) classify the domain on which the two partial members
have a common value. Thus
\[
 e_X(p,p)=d_X(p),\quad e_X(p,q)=e_X(q,p),\quad
 e_X(p,q)\le d_X(p)\wedge d_X(q),\quad
 e_X(p,q)\wedge e_X(q,r)\le e_X(p,r).
\]
Define in \(\mathcal E_0\)
\[
 \widehat d_X=h\,u_*d_X,\qquad \widehat e_X=h\,u_*e_X,
 \qquad A_X=\{p\in u_*\partial X:\widehat d_X(p)=\mathsf{true}_0\}.
\]
The finite-product identifications of \(u_*\) are understood.
On \(A_X\) let \(R_X\) be classified by \(\widehat e_X\),
and put \(H_hX=A_X/R_X\), using the established effective quotient
construction in \(\mathcal E_0\). The hats classify newly defined
predicates, not the images under \(u_*\) of the original subobjects.
Every formula is internal in the base and applies at all generalized
parameter objects.
\end{hbdefinition}

\begin{hbtheorem}{thm:p10-lawvere-closed-base-partials}{proved=annex:p10.lawvere-closed-base.partials}{The quotient functor and its exact finite-limit behavior}
The relation \(R_X\) is an equivalence relation. Forward image of
partial members defines an actual functor
\(H_h:\mathcal E\to\mathcal E_0\), with a natural singleton map
\(\eta_X:u_*X\to H_hX\). The quotient has the universal property
that maps \(H_hX\to Z\) correspond uniquely to maps
\(A_X\to Z\) equal on \(R_X\).
There is also an actual map \(\tau:H_hT\to T_0\) with
\(\tau\eta_T=h\); it is not asserted a classifier isomorphism.

The functor preserves terminal objects, products and equalizers, with
explicit inverse product maps and equalizer lifts obtained by restricting
partial members. It inverts every mono in the original class
\(\mathcal W_h\), not merely principal or finite-filter examples.
Consequently the exact ordinary localization property yields
\[
 \overline H_h:\mathcal E_h\longrightarrow\mathcal E_0,
 \qquad
 \overline H_h[U,f]=H_h(f)H_h(i_{U,X})^{-1}.
\]
This extension preserves finite limits of arbitrary localized arrows.
It is a new base functor, not a claimed factorization of \(u_*\).
\end{hbtheorem}

\begin{hbtheorem}{thm:p10-lawvere-closed-base-closed}{proved=annex:p10.lawvere-closed-base.closed}{Actual hom objects, composition and the closed comparison}
The cartesian closed structures of \(\mathcal E_h\) and
\(\mathcal E_0\) give the unit
\(1_0\cong\overline H_h1\) and natural comparison
\[
 \kappa_{X,Y}:\overline H_h(Y^X)
                  \longrightarrow(\overline H_hY)^{\overline H_hX},
\]
obtained by transposing transported evaluation. They satisfy all three
closed-functor compatibilities. In particular
\(\mathcal C_h(X,Y)=\overline H_h(Y^X)\) are actual base hom
objects, with identities transported from names of identity maps and
composition transported from internal composition. Associativity, unit
laws and the compatibility of \(\kappa\) with composition are proved
as equations of maps in \(\mathcal E_0\).
The closed comparison is not asserted an isomorphism or normal. The
ordinary maps of \(\mathcal E_h\) have natural names in these hom
objects; the naming maps need not be bijections onto their global elements.
The quotient and localization universal properties, and the exponential
transpose property fixing this product-induced comparison, are explicit
in the annex. No adjoint or geometric direct image is assumed for
\(\overline H_h\).
\end{hbtheorem}

\begin{hblemma}{lem:p10-lawvere-closed-base-countercases}{proved=annex:p10.lawvere-closed-base.countercases}{Two different base requirements that cannot be conflated}
For product \(u_*:\mathsf{Set}\times\mathsf{Set}\to\mathsf{Set}\)
and first-coordinate truth homomorphism, the new \(H_h\) is first
projection, whereas \(u_*\) cannot descend unchanged.
For diagonal \(u_*:\mathsf{Set}\to\mathsf{Set}\times\mathsf{Set}\)
and the identity truth homomorphism, \(\mathcal E_h=\mathsf{Set}\)
and \(H_h\) is diagonal. Its base hom object is
\((Y^X,Y^X)\); its global elements are pairs of functions, while
ordinary localized maps give only equal pairs. Thus even an invertible
closed comparison does not make the enriched global underlying category
the specified ordinary localization. Neither stronger requirement is
inserted into the source theorem.
\end{hblemma}

\begin{hbtheorem}{thm:p10-lawvere-closed-base-collapse}{proved=annex:p10.lawvere-closed-base.collapse}{The nondegenerate Boolean-valued two-valued collapse}
Let \(p:\mathcal E\to\mathcal S\) be a \emph{nondegenerate}
source Boolean-valued model over the abstract-set base. Then its internally
complete Boolean algebra \(B=p_*T\) admits a Boolean homomorphism
\(h:B\to2=T_{\mathcal S}\). This need not preserve arbitrary joins.
For this original datum \((p_*,h)\), the exact localization
\(\mathcal E\to\mathcal E_h\) is logical and its target is
nondegenerate, classical, two-valued, and still satisfies choice.
It has the constructed closed functor to \(\mathcal S\); it is
not claimed geometrically defined over \(\mathcal S\). In this
particular collapse case the naming maps are bijective, giving the
normal element-functor identification up to canonical isomorphism;
this uses the geometric original \(p_*\) and the proved choice
in \(\mathcal S\), unlike the unrestricted full-datum case.

The annex constructs \(h\) by bounded well-ordered proper-filter
histories, rather than assuming an ultrafilter lemma or a ZFC universe.
Nondegeneracy is necessary: the one-element Boolean algebra admits no
Boolean homomorphism to the nondegenerate algebra \(2\). The general
full-datum construction above includes the improper-filter/degenerate-base
boundary, but that boundary is not a nondegenerate two-valued model.
\end{hbtheorem}

These are independent editorial proofs. The original closed-base and
collapse statements remain historically announcements on p.~332. They
do not settle the unprinted normal-relative convention by inventing an
element-functor identity. They also do not supply the outstanding suitable
ccc/CH comparison and transfer, spectrum/locality/universality, or the
independent recursion-lifting/free-algebra arguments.
